\documentclass[a4paper, 11pt, french]{report}
\usepackage[utf8]{inputenc}	% package désignant l'encodage
\usepackage[T1]{fontenc}	% je sais pas à quoi il sert
\usepackage[francais]{babel} % indique qu'on écrit en français
\usepackage{layout}	% Permet de gérer les marges
\usepackage[top=2cm, bottom=2cm, left=2cm, right=2cm]{geometry}	% Permet de définir les marges
\usepackage{graphicx}	% Insertion des images
\usepackage{setspace}
\usepackage{titlesec}
\titleformat{\chapter}[hang]{\bf\huge}{\thechapter}{2pc}{}
\usepackage{fancyhdr}	% Pour gérer les en-tête / pied de page
\usepackage{placeins}	% Pour utiliser \FloatBarrier qui annule les flottants
\usepackage{listings}	
\usepackage{hyperref}


\graphicspath{{img/}{../}}	% Définir l'emplacement des images

\makeatletter

\def\clap#1{\hbox to 0pt{\hss #1\hss}}%
\def\ligne#1
{%
\hbox to \hsize{%
\vbox{\centering #1}}
}%

\def\haut#1#2#3
{%
\hbox to \hsize{%
\rlap{\vtop{\raggedright #1}}%
\hss
\clap{\vtop{\centering #2}}%
\hss
\llap{\vtop{\raggedleft #3}}}
}%
\def\bas#1#2#3
{%
\hbox to \hsize{%
\rlap{\vbox{\raggedright #1}}%
\hss
\clap{\vbox{\centering #2}}%
\hss
\llap{\vbox{\raggedleft #3}}}
}%
\def\logo#1#2#3
{%
\hbox to \hsize{%
\rlap{\vbox{\raggedright #1}}%
\hss
\clap{\vbox{\centering #2}}%
\hss
\llap{\vbox{\raggedleft #3}}}
}%
\def\maketitle
{%
\thispagestyle{empty}\vbox to \vsize
{%
\begin{tabular}{lcr}
\hbox
{
\mbox{%
\hspace{4pt}%
%\includegraphics[scale=0.8]{Images/IUT.png}%}
\hspace{20pt}
}
}%
&
\hspace{5cm}
&
\raggedright
\hbox
{
\mbox{%
\hspace{20pt}%
%\includegraphics[scale=0.6]{Images/synergie.jpg}%
\hspace{4pt}
}
}%
\end{tabular}

\vspace{1cm}
\haut{}{\@blurb}{}

\vfill
%\vspace{1cm}
\begin{flushleft}
\huge \@title
\end{flushleft}
\par
\hrule height 2pt
\par
\begin{flushright}
\Large \@author
\par
\end{flushright}
%\vspace{1cm}
\vfill
\hbox
{
\mbox{%
\hspace{4pt}%
\logo{}{
\includegraphics[scale=1]{univ_nantes.jpg}%
}{}
\hspace{4pt}
}
}%
\vfill
\bas{}{\@location, le \@date}{}
}%
\cleardoublepage
}
\def\date#1{\def\@date{#1}}
\def\author#1{\def\@author{#1}}
\def\title#1{\def\@title{#1}}
\def\location#1{\def\@location{#1}}
\def\blurb#1{\def\@blurb{#1}}
\date{\today}
\author{}
%\title{Rapport de stage de deuxième année de DUT Informatique}
\location{Nantes}\blurb{}
\makeatother
\title{{\sl\LARGE\textbf{TP Compilation - Langage de génération d’images SVG}}\normalsize}
\author{Guillaume Charon\\Jérôme Pages}
\blurb{%
Université de Nantes\\
UFR des Sciences et Techniques\\
Master ALMA 1
}


\begin{document}

	\pagestyle{plain}

	\maketitle	% Création de la page de garde
	\pagenumbering{arabic}

	\thispagestyle{empty}


	\renewcommand{\contentsname}{Sommaire} % Renomme la table des matières
	\setcounter{tocdepth}{2} % Définir la profondeur de la table des matières
	\addtocontents{toc}{%
	\protect\setlength{\baselineskip}{1.5em}	% Définir le saut de ligne dans la TDM
	\protect\setlength{\parskip}{0pt}	% Définir le saut de paragraphe dans la TDM
	}
	\tableofcontents\thispagestyle{plain}	% Table des matières
	%\pagenumbering{arabic}

	\begin{onehalfspace}
		%inclure le contenu
		% Work in progress

\chapter*{Introduction}

\chapter{Analyse du sujet}

	\section{Premières idées}
	Syntaxe
	
	\section{Mise en place de l'architecture}
	Structure de données utilisée pour stocker le xml 
	Construction du xml
	
	%TODO : parler de la génération qui se faisait à l'envers
	
	
\chapter{Développement du compilateur}
	
	\section{Partie 1 : La grammaire de base}

		\subsection{Construction de l'image}

		\subsection{Reconnaissance des formes basiques}

		\subsection{Ajout des attributs optionnels}
		% couleurs des formes, taille du texte, police d'écriture, etc...
	
		\subsection{Commentaires de fin de ligne}
	
		\subsection{Reconnaissance des formes complexes}
	
	\section{Partie 2 : Les instructions programmables}

		\subsection{Définition de variables et affectation}
		
		\subsection{Les opérations}

		\subsection{Les boucles}

\chapter*{Conclusion}



	\end{onehalfspace}

\end{document}
